\par\Needspace{6\baselineskip}\medskip
\begin{center}\hypertarget{native-185-article138}{}\pdfbookmark[2]{138.}{article-138}138.\end{center}
Negotium \emph{tertium} iam in eo consistit, ut duo loci geocentrici extremi corporis
coelestis, i.\,e. puncta $B,B''$, per circulum maximum iungantur, huiusque
intersectio cum circulo maximo $A'B'$ determinetur. Sit $B^*$ haec intersectio,
atque $\delta'-\sigma$ eius distantia a puncto $A'$, nec non $\alpha^*$ eius longitudo, $\beta^*$
latitudo. Habemus itaque, propterea quod $B,B^*,B''$ in eodem circulo maximo
iacent, aequationem satis notam
\[
\begin{aligned}
0={}&\operatorname{tang}\beta\sin(\alpha''-\alpha^*)-\operatorname{tang}\beta^*\sin(\alpha''-\alpha)\\
&+\operatorname{tang}\beta''\sin(\alpha^*-\alpha),
\end{aligned}
\]
quae, substituendo $\operatorname{tang}\gamma'\sin(\alpha^*-l')$ pro $\operatorname{tang}\beta^*$, sequentem formam induit
\[
\begin{aligned}
0={}&\cos(\alpha^*-l')\{\operatorname{tang}\beta\sin(\alpha''-l')-\operatorname{tang}\beta''\sin(\alpha-l')\}\\
&-\sin(\alpha^*-l')\{\operatorname{tang}\beta\cos(\alpha''-l')\\
&\qquad+\operatorname{tang}\gamma'\sin(\alpha''-\alpha)-\operatorname{tang}\beta''\cos(\alpha-l')\}.
\end{aligned}
\]
Quare quum sit $\operatorname{tang}(\alpha^*-l')=\cos\gamma'\operatorname{tang}(\delta'-\sigma)$, habebimus
\[
\begin{gathered}
\operatorname{tang}(\delta'-\sigma)=\\[.4ex]
\frac{\operatorname{tang}\beta\sin(\alpha''-l')-\operatorname{tang}\beta''\sin(\alpha-l')}{\begin{gathered}\cos\gamma'(\operatorname{tang}\beta\cos(\alpha''-l')-\operatorname{tang}\beta''\cos(\alpha-l'))\\+\sin\gamma'\sin(\alpha''-\alpha)\end{gathered}}\mathpunct{\cdot}
\end{gathered}
\]

Hinc derivantur formulae sequentes, ad calculum numericum magis accommodatae. Statuatur
\[
\begin{aligned}
7)\quad&\operatorname{tang}\beta\sin(\alpha''-l')-\operatorname{tang}\beta''\sin(\alpha-l')=S\\[.5ex]
8)\quad&\operatorname{tang}\beta\cos(\alpha''-l')-\operatorname{tang}\beta''\cos(\alpha-l')=T\sin t\\[.5ex]
9)\quad&\sin(\alpha''-\alpha)=T\cos t
\end{aligned}
\]
(art. 14, II), eritque
\[
10)\quad\operatorname{tang}(\delta'-\sigma)=\frac{S}{T\sin(t+\gamma')}\mathpunct{\cdot}
\]
Ambiguitas in determinatione arcus $\delta'-\sigma$ per tangentem inde oritur, quod circuli
maximi $A'B',BB''$ in \emph{duobus} punctis se intersecant: nos pro $B^*$ semper
\hypertarget{native-186}{}adoptabimus intersectionem puncto $B'$ proximam, ita ut $\sigma$ semper cadat inter
limites $-90^0$ et $+90^0$, unde ambiguitas illa tollitur. Plerumque tunc valor
arcus $\sigma$ (qui pendet a \emph{curvatura} motus geocentrici) quantitas satis modica erit,
et quidem generaliter loquendo secundi ordinis, si temporum intervalla tamquam
quantitates primi ordinis spectantur.

Quaenam modificationes calculo applicandae sint, si pro ecliptica aequator
tamquam planum fundamentale electum est, ex annotatione art. praec. sponte
patebit.

Ceterum manifestum est, situm puncti $B^*$ indeterminatum manere, si circuli
$BB'',A'B'$ omnino coinciderent: hunc casum, ubi quatuor puncta $A',B,B',B''$
in eodem circulo maximo iacerent, a disquisitione nostra excludimus.
Conveniet autem in eligendis observationibus eum quoque casum evitare, ubi
situs horum quatuor punctorum a circulo maximo parum distat: tunc enim situs
puncti $B^*$, qui in operationibus sequentibus magni momenti est, per
levissimos observationum errores nimis afficeretur, nec praecisione necessaria
determinari posset. --- Perinde punctum $B^*$ indeterminatum manere patet, quoties
puncta $B,B''$ in unum coincidunt\footnote{Sive etiam quoties sibi opposita sunt, sed de hoc casu non
loquimur, quum methodus nostra ad observationes tantum intervallum
complectentes non sit extendenda.}, in quo casu ipsius circuli $BB''$ positio
indeterminata fieret. Quamobrem hunc quoque casum excludemus, quemadmodum,
per rationes praecedentibus similes, talibus quoque observationibus
abstinendum erit, ubi locus geocentricus primus et ultimus in puncta sphaerae
sibi proxima cadunt.
